\chapter[Writing Began with More Than Poems]{The First Things Written Down Were Not Necessarily Poems}
\section{A Parcel Label}
Pick something up that probably lies at home: a parcel label.

The sheet stuck to a box, often casually torn off at collection, gives recipient, sender, address, telephone, weight, shipping charge, tracking number, barcode. Nobody reads it into tears. Neither novel, poem, nor social-media caption, it exists entirely to move an object accurately between people and make every intermediate stage answerable: where is it, who handled it, who bears responsibility?

Try a thought experiment. Five thousand years hence, after several civilisational turnovers, an archaeologist excavates this label, its lettering saved by plastic lamination. How might they describe our age?

Probably not as poetic. They would see vast logistics, millions of daily shipments, precise numbering and accounting at every stage, specialists reading and entering symbols. In one label: accounts, transport, responsibility, management.

This is no joke. What future archaeologists would do to us is precisely what today's do with earliest writing. Their discoveries differ greatly from imagined origins.

Many assume stories, songs, prayers, and overflowing hearts led humans to invent writing. Moving, but almost certainly wrong. The earliest generally accepted systematic written objects, Mesopotamian tablets around 5,300 years old, almost uniformly record bundles of reeds, beer jars, sheep, rations, recipients, outstanding balances.

Not poems. Bills.

That disappointing yet fascinating fact is our subject: earliest writing may concern barley in a store rather than poetry. Across five millennia from this start, writing changed power, literacy, and whose stories survive, occasionally letting an ordinary businessman's complaint outlast emperors' monuments.

\section{A Warehouse Five Thousand Years Ago}
Enter a scene.

Time: around 3200 BCE. Place: Uruk, a great Euphrates city in today's southern Iraq, among the world's largest settlements, perhaps tens of thousands strong. By contemporary standards, unmistakably a megacity.

More precisely, a store attached to the great temple. Mudbrick walls, timber beams and reed mats overhead, high-window light slanting through dust, grain, fermented beer smells. Reed baskets of dry barley stand on the floor; oil and beer jars line walls. Farmers drive in sheep while porters queue for daily rations.

A seated figure has wet tablets on a low stool. A trimmed reed with a small triangular tip presses successive wedge-shaped marks into clay, giving later cuneiform its name: Latin cuneus means wedge.

A sheep arrives: several impressions. Beer leaves: several more. Palm-wide tablets fill and dry hard in sunshine. Important ones go into a kiln, becoming nearly stone-durable ceramics. That later consequence is enormous foreshadowing.

This figure is a scribe, a reading-and-writing specialist, neither cultivator, fisher, nor potter. His entire work converts movements of goods into symbols. Uruk holds thousands of such tablets. Most excavated early examples are lists, rations, livestock registers, land and labour accounts. One famous tablet records workers' beer rations, a beer sign following every name: a five-thousand-year-old payslip paid in drink.

Stand there and notice sensory details. His hands press signs as fluently as an abacus user calculates. Waiting people crane toward incomprehensible clay: they recognise sheep, not writing. A supervisor periodically asks something; the scribe checks drying tablets and announces a number nobody in the queue can verify.

Remember the picture: many people around one tablet, only one able to read it. Everything following grows from here.

\section{For Whom Was Writing Invented?}
Set out the conflict before concluding.

Two accounts differ. Call the first expressive: writing grows from hearts full of stories, prayers, remembered dead, feelings demanding release. Its popularity suits romantic civilisation: beginnings should feature a poet or priest.

Call the second administrative: bigger cities, more goods, insufficient memory. Grain paid, rations received, land owned cannot remain oral without disputes. Symbols first settle accounts. Strong evidence supports it: earliest Mesopotamian tablets and Egyptian labels overwhelmingly concern administration and materials. Bone and ivory tags from an early Abydos royal tomb identify tribute's origins and quantities, doing your parcel label's job.

Material evidence favours accounting. Yet do not close the file: a deeper question, ours rather than archaeology's, appears:

\textbf{If writing began as management, did poetry, law, and love letters reform it, or did it never lose its instrumental nature?}

Is writing neutral like a knife, cutting food or people according to its holder? Or intrinsically biased toward registration, taxation, commands, and against preserving a farmer's feelings about weather, thereby automatically advantaging its owners?

This is not scholastic trivia. It shapes your reading of five millennia and today's forms, systems, platforms, all descendants of tablets and scribes.

Keep four things separate throughout. What happened, early surviving writing as accounts, is relatively undisputed evidence. Why writing arose in these places and times is competing explanation. How scribes understood their work belongs to their accounts. Whether writing was good or bad is evaluation; never mix it with the others.

Warm up your judgement: had writing first recorded poetry, how would the earliest surviving objects differ from actual finds?

\section{Who Reads and Writes, and Who Cannot?}
The crowded tablet with one reader holds our tension. Hear several voices.

\textbf{The scribes' voice.} They were proud professionals, not humble typists. Ancient Egypt's teaching text later called The Satire of the Trades mocks occupations successively: metalworkers' crocodile-skin fingers, mud-covered potters like rooting pigs, fishers among crocodiles, exhausted farmers unable to pay rent. Scribes carry no loads, avoid sun, receive respect, and issue orders. Copied by pupils, the text simultaneously teaches writing and superiority through literacy.

Schools left traces too. Mesopotamian scribal schools, Sumerian edubba, "tablet house", preserve many exercises. Archaeologists distinguish teachers' elegant models from awkward pupil copies. One Sumerian school story describes a child's day: beaten for lateness, mistakes, untidy clothing, then persuading his father to entertain and gift the teacher, improving treatment. Pupils suffered homework punishment five millennia ago. Some traditions run very deep.

\textbf{Give monopoly its full defence.} A natural response calls this knowledge monopoly: literate few bullying illiterate many. Not wrong, but hear the strongest defence of literacy's barriers before judging; its advocates are not all villains.

First, early scripts truly were difficult. Sumerian and Akkadian cuneiform had hundreds of signs, often multiple readings and meanings, as did Chinese characters and Egyptian hieroglyphs. Administrative competence required years from childhood. Technical difficulty, not only deliberate exclusion, raised barriers. Second, professional recording has public value: trained, accountable practitioners with standards provide more reliable accounts than casual universal scribbling, like today's accountants and contract-drafting lawyers. Third, records also protect ordinary people. Powerful parties can deny oral agreements; a written contract gives weaker parties at least paper or clay. Ordinary people would actively pay scribes, as we will see.

Yet intended or not, barriers yield power. Readers and writers occupy the intersection of goods, land, law, memory. Invisible hands can alter numbers, omit debts, decide whose complaint enters the record. Monopoly describes structure as well as motive: even universally virtuous scribes would occupy valuable positions.

\textbf{Beyond the counting room: writing leaves home.} Stopping here leaves mere bookkeeping. But invented tools acquire unforeseen uses. Within several centuries, writing gradually left warehouses:

\begin{itemize}
\item For law. Most famously, around 1754 BCE Hammurabi inscribed nearly three hundred rulings on a black stele over two metres tall and displayed it publicly.
\item For prayer and ritual. Hymns and funerary spells entered Egyptian tombs and Mesopotamian temple libraries.
\item For stories. Heroic traditions became epics, especially Gilgamesh, whose opening praises one who saw the deep and travelled the earth. Its narrative device imagines the hero engraving his experiences himself: storytellers borrow inscription's authority too.
\item For private life. Ordinary people hire scribes for letters, contracts, petitions. Next comes perhaps the oldest surviving consumer complaint.
\end{itemize}
Two other civilisations suggest early non-poetic writing is nearly global, though differently expressed. China's earliest discovered systematic characters are late Shang oracle bones, about 3,300 years old, on turtle shells and animal bones. Royal divination archives ask about harvest, war, childbirth, ancestral satisfaction. Fixed formulas date divination and identify the diviner asking whether the coming ten days bring harm. Archives, sacrifice, royal business, not poems; the Book of Songs follows centuries later. In Mesoamerica, Maya writing on monuments and altars records royal accessions, victories, calendrical rites: kingship's résumé rather than warehouse bills. Different contents share control by rulers and professional literates; none begins with folk songs.

\section[Thinking Bridge: Four Source Questions]{A Thinking Bridge: Four Questions for Ancient Writing}
Tablets, oracle bones, and stelae invite two extremes: total trust in written words or total dismissal as rulers' propaganda. Historians have a more useful portable tool, four questions:

\textbf{First: Who wrote it?} Royal secretary, temple priest, tax scribe, merchant hiring a writer? Identity shapes viewpoint and interests. Royal monuments omit humiliating defeats as résumés omit dismissals.

\textbf{Second: For whom?} Intended audiences determine what can and must be said. A public law stele addresses subjects and posterity, so its preface insists on justice. A merchant's private letter addresses one person and may disclose unvarnished anger and calculation.

\textbf{Third: Why?} Every inscription serves an immediate purpose: collecting debts or taxes, asserting authority, contacting gods, remembering. Purpose identifies strengths and inevitable omissions. A store list faithfully counts thirty sheep and stays silent about the driving farmer's thoughts, not necessarily concealing them but lacking such a field.

\textbf{Fourth: How did it survive?} Easily overlooked, especially sharp. We do not see everything ancient people wrote, only writing on sufficiently durable material, in suitable places, fortuitously excavated and deciphered. Tiny remnants survive successive filters. Statisticians call treating these as representative \textbf{survivorship bias}. Fire-, water-, time-resistant clay and stone favour states, temples, wealth. Poor people lived too; their voices occupied perishable things or went unwritten.

Apply all four again. Parcel label: logistics company, sorting system, tracking, hypothetical survival through plastic. Oracle inscription: royal diviner, ancestors and court, divination archive, survival buried in pits. Hammurabi: royal commission, subjects and gods, legitimate rule, giant-stone preservation.

One conclusion emerges: \textbf{hard surfaces record not necessarily what mattered most, but what people had power to write and what happened to survive}. Put it in your notebook. It applies to historical trivia feeds as well as archaeology.

\section{A Tablet Full of Anger}
Read primary material armed with those questions.

An eighteenth-century BCE tablet excavated at Ur, now in the British Museum, carries Nanni's letter to copper merchant Ea-nasir. Nanni is neither king nor priest, just a businessman. Translations differ in wording, so here is its sense: you took payment and promised fine ingots but supplied inferior copper; my servants crossed hostile territory to recover money, yet you sent them away empty-handed and insulted them; who do you think I am, treating me so? Henceforth I accept only fine copper and reject the rest.

Who writes? An ordinary paying, aggrieved merchant, probably dictating to a scribe. For whom? The copper merchant alone. Why? Recovery, explanation, future bargaining position. Survival? A tablet in an abandoned Ur house outlasted nearly four millennia because clay had hardened through firing.

These questions explain its repeatedly exhibited value: \textbf{nobody intended it for posterity}. Kings build monuments for ages; Nanni only wanted that dishonest trader to repay him. Precisely its lack of commemorative intention leaves real anger, negotiation, and social relationships unpolished, accidentally captured by writing. Without it, his grievance vanishes; with it, an unknown person's complaint outlasts many royal names. Calling it history's oldest bad review is amusing and meaningful. Writing's extension of rule is familiar; its preservation of ordinary voices often happens casually, accidentally, unpredictably.

Contrast Hammurabi's preface, claiming gods granted power to prevent strong oppressing weak and defend orphans and widows before listing laws. Four questions: royal order; subjects, gods, posterity; legitimation; durable black stone surviving nearly intact millennia. This need not mean falsehood: proclaimed and actual justice can be tested separately. It shows writing's two directions: deliberately designed top-down commemoration versus accidental bottom-up preservation. Historians use both, never reading them identically.

\section{Why Writing Appeared: Three Accounts}
Return to invention's causes, setting genuinely different explanations, strengths, and limits beside one another.

\textbf{Explanation 1: Accounting: administrative demand.} Strongest in Mesopotamia. Denise Schmandt-Besserat proposed an influential sequence: tokens counted sheep and sacks millennia before writing; goods in transit used hollow clay envelopes as seals; surface impressions displayed contents without breaking them; eventually impressions replaced enclosed tokens, becoming tablet symbols. Elegant and supported by material sequences, this matches earliest account texts. Yet its regional chain remains disputed in places and cannot explain Chinese divination or Maya royal chronology, which lack that token sequence.

\textbf{Explanation 2: Gods and authority: ritual power.} Better fitting oracle bones and Maya monuments, this makes earliest writing ancestral communication, royal declaration, ritual record, sacred and political from inception rather than economic. It recognises writing's entanglement with faith and power, but struggles with Mesopotamia's mountains of beer rations. Surely every beer distribution cannot be spirit communication.

\textbf{Explanation 3: The question is mistaken: multiple origins and causes.} A singular invention question presumes one worldwide answer. Writing arose independently at least in Mesopotamia, China, and Mesoamerica, with Egypt's independence debated. One needed warehouses, another ancestors, another royal biographies. Better: when social complexity exceeds oral memory's capacity for something, accounts or rites, external memory becomes necessary; the particular something varies. Most inclusive, least sharp: accommodating everything makes falsification and bold prediction difficult.

Notice that causes are disputed while earliest textual types form shared evidence. Good historical arguments agree on a factual ledger and beat different explanatory drums. Next time scholars quarrel, locate their disputed layer.

\section[Further Challenge: Writing and Thought]{A Deeper Challenge: Does Writing Change Thought Itself?}
So far, writing changes society's power, records, memory. Does it change \textbf{thought itself}? British anthropologist Jack Goody helped lead this twentieth-century dispute.

In The Domestication of the Savage Mind (1977) and other works, Goody argues writing creates mental forms \textbf{impossible in purely oral society}, not merely transferring speech to clay. Lists are exemplary. Spoken sheep, oil, beer quantities vanish with sound; written names and numbers become columns. A table permits instant comparison, grouping, spotting empty cells, reordering. Classification, comparison, sorting, exceptions, seemingly natural thought, are partly trained by written tables. Walter Ong's Orality and Literacy (1982) pushes further: oral thought is contextual, additive and formulaic rather than analytically dissecting, remembered through narrative patterns; writing makes thought abstract and analytical. Writing reorganises thinking, not merely records it.

Powerful, but a dangerous reading demands direct rejection: literate civilisations smarter, nonliterate peoples mentally inferior. Specifically wrong. Later critics argue Goody makes literacy's role too linear. Its effects depend on use: universal literacy confined to scripture recitation need not beat limited literacy with vigorous debate in analytical thought. Schools, exams, argument, publishing, the surrounding institutions rather than writing alone, shape thinking.

Fix the lesson with an easily misread \textbf{counterexample: the Inca Empire}. Fifteenth-century Andean rulers governed over ten million, built twenty or thirty thousand kilometres of roads, managed mountain-to-coast granaries, and levied population-based labour and material taxes, without writing. They used quipu, coloured knotted cords encoding numbers and perhaps categories through material, colour, knot positions and spacing, maintained by specialists. Knots supported precise imperial tax administration. Equating no writing with no governance or history badly misreads the Inca.

Oral culture has positive strengths too. Homer's Iliad and Odyssey exceed twenty-seven thousand lines. Twentieth-century researcher Milman Parry demonstrated their composition and transmission through oral formulas outside writing, singers using repeated phrases such as rosy-fingered dawn for memory and improvisation. West African griots recite generations of royal descent, wars, alliances, giving courtroom and ceremonial testimony with archival weight. Calling them uncultured reveals not their deficiency but literate people's long monopoly over culture's definition.

The conclusion has two layers. Goody adds mental organs through writing: tables, lists, abstract categories. Inca and griots show their absence does not abolish thought, governance, history. \textbf{Literate and nonliterate societies both possess history; only its home differs: clay and paper, or memory, ritual, knots.} Memory can die with people, its weakness; tablets allow repeated checking and contradiction, their strength. Yet remember that checkable material remains only the little that happened to be written.

\section{Today: We All Live in the Warehouse}
This ancient question watches from your pocket.

Today's ordinary people are history's most recorded. An ancient farmer might leave one tax-list name; your medical birth record, school file, grades, payments, journeys, chats, searches could pile more tablets than all Uruk's scribes pressed in their lives. Writing and descendant databases never stopped expanding records; warehouses moved into servers.

New warehouses have new scribes. Literacy now also means reading data, writing code, understanding system labels. Credit scores, recommendations, risk models are new cuneiform: masses live around them, few specialists understand their construction. Ancient illiteracy upgrades rather than disappears. You recognise every app word while not knowing what is recorded or calculated about you.

But modern reversals deserve individual attention.

First, ordinary voices survive at unprecedented scale. Nanni unhappily hired a tablet scribe; anyone now posts reviews, messages, videos nearly free, potentially lasting long. The narrow opening for accidental survival becomes a square.

Second, survival's logic grows stranger. Clay relies on material durability; digital records on functioning platforms, supported formats, continued payments. Five-thousand-year fired clay remains legible, while twenty-year-old floppy disks and closed websites may be lost forever. Vast digital memory can have remarkably short shelf life. Future historians may confront not scarcity but evidence fragmented into inaccessible formats.

Third, speech revives online instead of dying. Ong observed secondary orality through radio, telephone, recording: voices return atop written and archival infrastructure. Voice messages, podcasts, video calls belong here. Once speech and writing followed one another; now they share your phone. What griots protected by recitation may survive through a recording forwarded a hundred thousand times.

Fourth, authority to write changes costume. Ancient scribes chose records; designers choose fields and abnormal behaviours. Field names are inconspicuous modern power: whether a form records rejection reasons affects appeals; whether attendance logs overtime reasons affects recognition of effort. Hammurabi publicly inscribed laws; many rules governing you hide in unseen code. No programming skill is needed to notice. Next form, consent button, or system score, remember Uruk's one reader among many, then ask who wrote this, for whom, why, and how long it will last.

\section{Your Turn: Leave a List, or a Poem?}
Return to the parcel label and Nanni's angry tablet.

Suppose archaeologists five millennia hence can receive exactly \textbf{one} written object from you to understand your life and age. What will you leave?

A list: timetable, monthly expenses, annual steps? Lists do not lie; more honestly than self-introductions they reveal time, money, bodily effort, as Uruk's rations explain its workings better than hymns.

A letter like Nanni's, addressed to a specific person with specific grievances, requests, temperament? It carries a relationship's warmth and living voice.

Or poem, or story? Writing's ventures away from home, unforeseen by its accounting inventors: tools of management holding what management cannot contain.

Before choosing, balance our lessons: writing need not preserve what matters most; survivors favour empowered writers and durable materials; lists are more honest than poems, poems longer-lived than lists; a civilisation's self-introduction is a survivorship-filtered résumé. Remember unwritten Inca administration and griot genealogies: not writing is also an answer.

Choose with reasons. No standard answer exists, but every choice determines what someone five thousand years later can see and can never see. Every person taking up a pen has done this since that ancient warehouse.
